\section{Gowers-Hatami for classical BLR}

In classical BLR, the mapping is from a field $\mathbb{F}_2^n$. 
First let us relate the BLR test to the notion of approximate representation.

\begin{proposition}[BLR test induces an approximate representation]\label{lem:blr-approx-rep}
Let $f:\mathbb{Z}_2^n\to \mathbb{Z}_2$ satisfy
\[
\Pr_{x,y}[f(x)+f(y)=f(x+y)] \ge 1-\epsilon.
\]
Define $F:\mathbb{Z}_2^n\to \{\pm1\}$ by $F(x)=(-1)^{f(x)}$. Then
\[
\mathbb{E}_{x,y}\big|F(x)F(y)-F(x+y)\big|^2 \le 4\epsilon,
\]
so $F$ is an $(2\epsilon,1)$-approximate representation.
\end{proposition}

\begin{proof}
    First observe that for all $x,y\in G$,
    \[
    F(x)F(y) = (-1)^{f(x)+f(y)}.
    \]
    Hence
    \[
    F(x)F(y)=F(x+y) \quad \Longleftrightarrow \quad f(x)+f(y)=f(x+y)\ (\mathrm{mod}\ 2).
    \]

    Define the indicator of BLR success:
    \[
    \mathbf{1}_{\mathrm{good}}(x,y) := \mathbf{1}[f(x)+f(y)=f(x+y)].
    \]
    Then
    \[
    \Pr_{x,y}[\mathbf{1}_{\mathrm{good}}=1] \ge 1-\epsilon.
    \]

    Now compute the squared error. Since $F(x)\in\{\pm1\}$,
    \[
    |F(x)F(y)-F(x+y)|^2 =
    \begin{cases}
    0 & \text{if } F(x)F(y)=F(x+y),\\
    4 & \text{otherwise}.
    \end{cases}
    \]

    Taking expectation over $(x,y)$,
    \begin{align*}
    \mathbb{E}_{x,y}|F(x)F(y)-F(x+y)|^2
    &= 4 \Pr_{x,y}[F(x)F(y)\neq F(x+y)] \\
    &= 4 \Pr_{x,y}[f(x)+f(y)\neq f(x+y)] \\
    &\le 4\epsilon.
    \end{align*}
\end{proof}

The Gowers-Hatami theorem can be applied to show that if a function $f: \mathbb{Z}_2^n \to \mathbb{F}_2$ passes the BLR linearity test with high probability, 
then there exists a linear function $\ell: \mathbb{F}_2^n \to \mathbb{F}_2$ that is close to $f$ in terms of Hamming distance $\delta(f, \ell)$ defined in Def.~\ref{def:relative-hamming-distance}.
To do this, we first need to apply the Gowers-Hatami theorem to the approximate representation $F(x)=(-1)^{f(x)}$ given by Proposition~\ref{lem:blr-approx-rep}.

\begin{corollary}[Gowers-Hatami Correlation for BLR]\label{cor:gh-Z2}
    Let $F: \mathbb{Z}_2^n \to U(1)$ be an $(\epsilon, \sigma)$ approximate representation. 
    Then there exists a unit vector $v$ and an exact representation $G$ such that 
    \be
    \mathbb{E}_{x \in \mathbb{Z}_2^n} \Re \langle F(x), \langle v, G(x) v\rangle \rangle \geq 1 - \epsilon.
    \ee
    
\end{corollary}

Next, we would like to conclude from the above corollary that there exists a linear function $\ell: \mathbb{Z}_2^n \to  \mathbb{Z}_2$ such that 
$\Pr_x[f(x) \neq \ell(x)] \leq \epsilon/2$. To do this, we need to use a fact that about the regular representation of $\mathbb{Z}_2^n$.


% need to prove this fact
\begin{fact}[Regular representation of $\mathbb{Z}_2^n$]
    Let $\rho(a)$ act on $L=\{f:\mathbb{Z}_2^n\to\mathbb{C}\}$ by $(\rho(a)f)(x)=f(x+a)$. 
    Then in the Fourier basis $\{\ket{S}\}$, we have
    \[
    \rho(a) = \sum_{S\subseteq[n]} \chi_S(a)\ket{S}\bra{S}.
    \]
\end{fact}

\begin{lemma}[Gowers-Hatami implies near linearity]\label{lem:gh-linearity}
    Let $F: \mathbb{Z}_2^n \to \mathbb{C}$ be an $(\epsilon, 1)$ approximate representation of $\mathbb{Z}_2^n$ from the BLR test 
    such that $F(x) = (-1)^{f(x)}$ for some function $f: \mathbb{Z}_2^n \to \mathbb{Z}_2$.
    Then there exists a linear function $\ell: \mathbb{F}_2^n \to \mathbb{F}_2$ such that the distance between $f$ and $\ell$ is bounded
    \be
    \delta(f, \ell) \leq \epsilon/2.
    \ee
\end{lemma}


\begin{proof}
    By Corollary~\ref{cor:gh-Z2}, there exists a representation $G$ and a unit vector $v$ such that
    \[
        \mathbb{E}_{x} \Re \langle F(x), \langle v, G(x)v\rangle \rangle \ge 1 - \epsilon.
    \]
    Using the diagonal form of $G(x)$ in the character basis,
    \[
        G(x) = \sum_S \chi_S(x) \ket{S}\bra{S}.
    \]

    The inner product condition can be rewritten as
    \[
       \sum_{S} p_S  \mathbb{E}_{x} \Re \langle F(x), \chi_S(x) \rangle \ge 1 - \epsilon,
    \]
        where $p_S := |\bra{v}S\rangle|^2$ and $\sum_S p_S = 1$.
    Therefore,
    \[
        \sum_S p_S \widehat{F}(S) \ge 1 - \epsilon,
    \]
    where $\widehat{F}(S):=\mathbb{E}_x[F(x)\chi_S(x)]$ is the real-valued Fourier coefficient of $F$ at $S$.

    Since $(p_S)_S$ is a probability distribution, there exists some $S^*$ such that
    \[
        \widehat{F}(S^*) \ge 1 - \epsilon.
    \]
    This means that there is a large Fourier coefficient at $S^*$.
    Define $\ell(x)=x\cdot S^*$. Then $\chi_{S^*}(x)=(-1)^{\ell(x)}$, and
    \[
        \mathbb{E}_x[F(x)\chi_{S^*}(x)] = 1 - 2\Pr_x[f(x)\neq \ell(x)].
    \]
    Hence
    \[
        1 - 2\delta(f,\ell) \ge 1 - \epsilon,
    \]
    which implies
    \[
        \delta(f,\ell) \le \epsilon/2.
    \]
\end{proof}

\begin{proposition}[Equivalence for BLR test]\label{prop:equiv_blr_z2}
    Let $ \Pr[V_{\text{BLR}}^f=0]$ denote the probability that the BLR verifier rejects a function $f: \mathbb{F}_2^n \to \mathbb{F}_2$.
    Let $\delta(f, \mathcal{LIN})$ denote the distance between $f$ and the closest linear function in $\mathcal{LIN}$.
    Then the following statements are equivalent:
    \begin{itemize}
        \item If $\Pr[V_{\text{BLR}}^f=0] \leq \epsilon$, then $\delta(f, \mathcal{LIN}) \leq \epsilon$.
        \item $\Pr[V_{\text{BLR}}^f=0] \geq \delta(f, \mathcal{LIN})$.
    \end{itemize}
\end{proposition}    


Now we provide the proof of Soundness of BLR test using the Gowers-Hatami theorem.

\begin{theorem}[Soundness of the $Z_2$ BLR]
\label{thm:z2-blr-soundness}
For every function $f:\mathbb{F}_2^n\to\mathbb{F}_2$,
\[
    \Pr[V_{\text{BLR}}^f=0] \ge \delta(f,\mathcal{LIN}).
\]
Equivalently,
\[
    \Pr[V_{\text{BLR}}^f=1] \le 1-\delta(f,\mathcal{LIN}).
\]
\end{theorem}

\begin{proof}
    Let $F: \mathbb{Z}_2^n \to \{-1,1\}$ be defined by $F(x)=(-1)^{f(x)}$, as given by Proposition~.\ref{lem:blr-approx-rep}.
    $F$ is an $(2\epsilon,1)$ approximate representation of $\mathbb{Z}_2^n$.
    Applying Gowers-Hatami theorem by Lemma.~\ref{lem:gh-linearity},
    then there exists a linear function $\ell$ such that $\delta(f, \ell) \leq \epsilon$ so that $\delta(f,\mathcal{LIN})\le \epsilon$.

    This proves the implication 
    $\Pr[V_{\text{BLR}}^f=0]\le \epsilon \Rightarrow \delta(f,\mathcal{LIN})\le \epsilon$,
    which is equivalent to
    \[
        \Pr[V_{\text{BLR}}^f=0]\ge \delta(f,\mathcal{LIN}),
    \]
     by Proposition~\ref{prop:equiv_blr_z2}.
\end{proof}